%------------------------------------------------------------------------------%
\begin{question}
  Dados
  \(
    \vec{u} = \vetor{2}{-1}
  \) e
  \(
    \vec{v} = \vetor{-4}{3}
  \)

  \bigskip
  Determine $\lambda$ de modo que
  \(
    \norm{2\vec{u} + \vec{v} + \vetor{\lambda}{1}} = 3
  \)
\end{question}

%------------------------------------------------------------------------------%
\begin{resolution}{Solução}
  \begin{align*}
    \onslide<+->{\vec{w}}
    \onslide<+->{& = 2\vec{u} + \vec{v} + \vetor{\lambda}{1}              \\}
    \onslide<+->{& = 2\vetor{2}{-1} + \vetor{-4}{3} + \vetor{\lambda}{1}  \\}
    \onslide<+->{& = \vetor{4}{-2} + \vetor{-4}{3} + \vetor{\lambda}{1}   \\}
    \onslide<+->{& = \vetor{4-4+\lambda}{-2+3+1}                          \\}
    \onslide<+->{& = \vetor{\lambda}{2}}
  \end{align*}
\end{resolution}

%------------------------------------------------------------------------------%
\begin{resolution}{Solução}
  \begin{align*}
    \onslide<+->{\norm{\vec{w}}}
    \onslide<+->{= \norm{\vetor{\lambda}{2}}}
    \onslide<+->{& = 3 \\}
    \onslide<+->{\sqrt{\lambda^2 + 2^2} & = 3        \\}
    \onslide<+->{\lambda^2 + 4          & = 3^2      \\}
    \onslide<+->{\lambda^2              & = 9 - 4    \\}
    \onslide<+->{\lambda^2              & = 5        \\}
    \onslide<+->{\abs{\lambda}          & = \sqrt{5} \\}
    \onslide<+->{\lambda = \sqrt{5}     & }
    \onslide<+->{\quad \text{ou} \quad \lambda = -\sqrt{5}}
  \end{align*}
\end{resolution}

%------------------------------------------------------------------------------%
\endinput
